% Point to the main file
\documentclass[../../Main.tex]{subfiles}

% ========== Start the document ==========
\begin{document}

\section{Section 2.6 - Matrices}

\includegraphics[width=15cm]{directed_graph.png}

\begin{definition}[Directed Graph]
    This is a \yellowbox{directed graph} and could be used to model:

    \begin{itemize}
        \item How webpages are linked
        \item Flights between cities
        \item Who knows who
    \end{itemize}

    How can a computer store this directed graph?

\end{definition}

\begin{definition}[Matrix]
    We will use a \yellowbox{matrix}, which is a rectangular array of numbers
\end{definition}

For this network, we define a matrix

$$ A = \left[ a_{ij} \right] $$

Where

\[
a_{ij} = 
    \begin{cases}
        0 & \text{if web page $i$ doesn't link to $j$} \\
        1 & \text{if web page $i$ does link to $j$} \\
    \end{cases}
\]

\begin{align*}
    i &= \text{row (starts at 1)} \\
    j &= \text{column (starts at 1)} \\
\end{align*}

The resulting matrix would look like

\[
\begin{bmatrix}
    a_{11} & a_{12} & a_{13} & a_{14} \\
    a_{21} & a_{22} & a_{23} & a_{24} \\
    a_{31} & a_{32} & a_{33} & a_{34} \\
    a_{41} & a_{42} & a_{43} & a_{44} \\
\end{bmatrix}
\]

For this network, the matrix is

\[
A = 
\begin{bmatrix}
    1 & 1 & 0 & 0 \\
    1 & 0 & 1 & 0 \\
    0 & 0 & 0 & 1 \\
    0 & 1 & 0 & 1 \\
\end{bmatrix}
\]

Let's modify the graph now

\includegraphics[width=15cm]{directed_graph_other.png}

This matrix is

\[
B = 
\begin{bmatrix}
    1 & 1 & 0 & 0 \\
    0 & 1 & 1 & 1 \\
    0 & 0 & 0 & 0 \\
    1 & 0 & 1 & 0 \\
\end{bmatrix}
\]

\subsubsection{Matrix Addition and Subtraction}

We define the addition of matrices by adding corresponding components

\[
A + B = 
\begin{bmatrix}
    2 & 2 & 0 & 0 \\
    1 & 1 & 2 & 1 \\
    0 & 0 & 0 & 1 \\
    1 & 1 & 1 & 0 \\
\end{bmatrix}
\]

The same can be said for Subtraction

\[
A - B = 
\begin{bmatrix}
    0 & 0 & 1 & 1 \\
    1 & -1 & 1 & -1 \\
    0 & 0 & 0 & 1 \\
    -1 & 1 & -1 & 1 \\
\end{bmatrix}
\]

\begin{note}
    Two matrices can only be added/subtracted if they're the same size
\end{note}

\begin{question}
    Going back to the first network, how many ways can we start at 1 and return to 1 after passing through
    exactly 1 webpage?
\end{question}

\begin{align*}
    1 &\rightarrow 1 \rightarrow 1 \\
    1 &\rightarrow 2 \rightarrow 1 \\
    a_{12} &\rightarrow a_{21} \\
\end{align*}

\[
\left[
    \begin{array}{cccc}
        \rowcolor{yellow} \overset{\text{$a_{11}$}}{1} & \overset{\text{$a_{12}$}}{1} & 0 & 0 \\
        1 & 0 & 1 & 0 \\
        0 & 0 & 0 & 1 \\
        0 & 1 & 0 & 1 \\
    \end{array}
\right]
\left[
    \begin{array}{>{\columncolor{green}}c ccc}
        \overset{\text{$a_{11}$}}{1} & 1 & 0 & 0 \\
        \overset{\text{$a_{21}$}}{1} & 0 & 1 & 0 \\
        0 & 0 & 0 & 1 \\
        0 & 1 & 0 & 1 \\
    \end{array}
\right]
=
\left[
    \begin{array}{cccc}
        2 & 1 & 1 & 0 \\
        1 & 1 & 0 & 1 \\
        0 & 1 & 0 & 1 \\
        1 & 1 & 1 & 1 \\
    \end{array}
\right]
\]

\begin{example}
    Let

    \[
    A =
    \begin{bmatrix}
        1 & 2 & 3 \\
        4 & 5 & 6 \\
    \end{bmatrix}
    , \;
    B =
    \begin{bmatrix}
    1 & 0 \\
    -1 & 1 \\
    0 & 1 \\
    \end{bmatrix}
    , \;
    C =
    \begin{bmatrix}
    1 & 2 \\
    2 & 4 \\
    \end{bmatrix}
    \]

    \begin{enumerate}
        \item
        {
            $A + B$

            \begin{solution}
                \textbf{Undefined}, due to dimensions
            \end{solution}
        }

        \item 
        {
            $6A$

            \begin{solution}
                \[
                    6A =
                    \begin{bmatrix}
                        6 & 12 & 18 \\
                        24 & 30 & 36 \\
                    \end{bmatrix}
                \]
            \end{solution}
        }

        \item
        {
            $AB$

            \begin{solution}
                \[
                AB =
                \left[
                    \begin{array}{ccc}
                        \rowcolor{yellow} 1 & 2 & 3 \\
                        4 & 5 & 6 \\
                    \end{array}
                \right]
                \left[
                    \begin{array}{>{\columncolor{green}}c ccc}
                        1 & 0 \\
                        -1 & 1 \\
                        0 & 1 \\
                    \end{array}
                \right]
                =
                \left[
                    \begin{array}{cc}
                        -1 & 5 \\
                        -1 & 11 \\
                    \end{array}
                \right]
                \]
            \end{solution}
        }

        \item
        {
            $BA$

            \begin{solution}
                \[
                BA =
                \left[
                    \begin{array}{cc}
                        \rowcolor{yellow} 1 & 0 \\
                        -1 & 1 \\
                        0 & 1 \\
                    \end{array}
                \right]
                \left[
                    \begin{array}{>{\columncolor{green}}c cc}
                        1 & 2 & 3 \\
                        4 & 5 & 6 \\
                    \end{array}
                \right]
                =
                \left[
                    \begin{array}{ccc}
                        1 & 2 & 3 \\
                        3 & 3 & 3 \\
                        4 & 5 & 6 \\
                    \end{array}
                \right]
                \]
            \end{solution} 
        }

        \item
        {
            $BC$

            \begin{solution}
                \[
                BC = 
                \left[
                    \begin{array}{cc}
                        \rowcolor{yellow} 1 & 0 \\
                        -1 & 1 \\
                        0 & 1 \\
                    \end{array}
                \right]
                \left[
                    \begin{array}{>{\columncolor{green}}c c}
                        1 & 2 \\
                        2 & 4 \\
                    \end{array}
                \right]
                =
                \left[
                    \begin{array}{cc}
                        1 & 2 \\
                        1 & 2 \\
                        2 & 4 \\
                    \end{array}
                \right]
                \]
            \end{solution}
        }

        \item
        {
            $CB$

            \begin{solution}
                \textbf{Undefined}, due to dimensions
            \end{solution}
        }

        \item
        {
            $C^2 = CC$

            \begin{align*}
            C^2 = CC
            =
            \left[
                \begin{array}{cc}
                    \rowcolor{yellow} 1 & 2 \\
                    2 & 4 \\
                \end{array}
            \right]
            \left[
                \begin{array}{>{\columncolor{green}}c c}
                    1 & 2 \\
                    2 & 4 \\
                \end{array}
            \right]
            &=
            \left[
                \begin{array}{cc}
                    5 & 10 \\
                    10 & 20 \\
                \end{array}
            \right] \\
            &= 5C
            \end{align*}
        }

        \item
        {
            $C^{11}$

            \begin{align*}
                C^2 &= 5C \\
                C^3 &= C^2 \cdot C \\
                &= 5C \cdot C \\
                &= 5C^2 \\
                &= 5^2 C \\
                C^{11} &= 5^{10} C = 5^{10} 
                \left[
                    \begin{array}{cc}
                        1 & 2 \\
                        2 & 4 \\
                    \end{array}
                \right]
            \end{align*}
        }
    \end{enumerate}
\end{example}

Okay let's now define some properties

\subsection{Properties of Matricies}

\begin{definition}[Additive Identity]
    We define the \yellowbox{Additive Identity} for matrices as

    \begin{align*}
        \underset{\text{$n \times m$}}{A} + \vec{0} = \underset{\text{$n \times m$}}{A}, 
        \;
        \vec{0} =
        \begin{bmatrix}
            0 & 0 & \cdots & 0 \\
            0 & 0 & \cdots & 0 \\
            \vdots & \vdots & \ddots & \vdots \\
            0 & 0 & \cdots & 0 \\
        \end{bmatrix}
    \end{align*}

\end{definition}

\begin{definition}[Multiplicative Identity]
    We define the \yellowbox{Multiplicative Identity for matrices as}

    \begin{align*}
        \underset{\text{$n \times m$}}{A} \cdot I_n = A,
        \;
        I_n =
        \begin{bmatrix}
            1 & 0 & \cdots & 0 \\
            0 & 1 & \cdots & 0 \\
            \vdots & \vdots & \ddots & \vdots \\
            0 & 0 & \cdots & 1 \\
        \end{bmatrix}
    \end{align*}

    But $A$ must be \textbf{square}!
\end{definition}

Some matrices have a multiplicative inverse such that

\begin{definition}[Multiplicative Inverse]
    We define the \yellowbox{Multiplicative Inverse} of matrices as

    $$ A \cdot B = I $$

    But only \textbf{some} matrices have them, not all do!

    Also, $A$ must be square to even qualify for possibly having one

    For this class, we define

    \begin{align*}
        A =
        \begin{bmatrix}
            a & b \\
            c & d \\
        \end{bmatrix}
        ,
        \;
        A^{-1} =
        \frac{1}{ad - bc}
        \begin{bmatrix}
            d & -b \\
            -c & a \\
        \end{bmatrix}
    \end{align*}
\end{definition}

\begin{example}
    We have
    
    \begin{align*}
        \begin{bmatrix}
            4 & 3 \\
            7 & 5 \\
        \end{bmatrix}^{-1}
        =
        \frac{1}{20 - 21}
        \begin{bmatrix}
            5 & -3 \\
            -7 & 4 \\
        \end{bmatrix}
    \end{align*}

\end{example}


% ========== End the document ==========
\end{document}